\documentclass[10pt, a4paper]{article}
% The next line loads some packages you will need
\usepackage[UKenglish]{datetime}
\usepackage{a4, graphicx, amsmath, amssymb}
\usepackage[vmargin=30mm,hmargin=42mm]{geometry}
\thispagestyle{empty} 
\begin{document}
\begin{flushright}
\newdateformat{UKvardate}{%
\THEDAY\ \monthname[\THEMONTH] \THEYEAR}
\UKvardate
{\em Himanshu Madan, 201501359, \today}
\end{flushright}
\begin{center}
    {\bf \Large MATH5747/5825 assignment 1: {\LaTeX} }
\end{center}

{\indent Recreate this page, changing only the name, ID, and date above. You must use the {\LaTeX} template which is provided in the module area on Minerva. Five items you must use are listed below. (If you are missing any of them, then you are probably doing something wrong!)}
\begin{enumerate}
    \item Lists, which can be of type
        \begin{itemize}
            \item 
                itemize
            \item 
                enumerate
            \item
                description
        \end{itemize}
    \begin{description}
        \item({\bf Note:} You must use all three on this page.)
    \end{description}
    \item 
        Tables, such as table \ref{tab:Table1}
    \begin{table}[h!]
    \begin{center}
        \begin{tabular}{c|c}
        {\em English} & {\em Greek}\\
        \hline
        a, A & $\alpha$\\
        b, B & $\beta$\\
        d, D & $\delta$, $\Delta$\\
        \vdots & \vdots\\
        w, W & $\omega$, $\Omega$
        \end{tabular}
    \end{center}
    \caption{English and Greek letters.}
    \label{tab:Table1}
    \end{table}
    \item
    The commands \verb|\label| and \verb|\ref|.
    \item 
    Figures, such as figure 1 below. The file you will need, \verb|Leeds_Logo.pdf|, is available in pdf format  from Minerva.
    \begin{figure}[h]
	\includegraphics[width=6.5cm]{Leeds_Logo}
	\centering
	\caption{The University of Leeds logo.}
    \end{figure}\\
    It has been scaled to be 6.5 cm wide by using the option \verb|width=6.5cm| to the \verb|\includegraphics| command.
    
    {\em R} graphics can be exported as pdf files. The method depends on your oper-\\ating system.
    \begin{description}
	    \item
	    [\bf Windows] Click on the graphics window and select \texttt{file} $\rightarrow$ \texttt{save as}  $\rightarrow$\\ \texttt{pdf}.
    \end{description}
    \begin{description}
	    \item[linux/UNIX] 
	        Use the {\em R} command \texttt{dev.copy2pdf(file="filename.pdf").}\\ (Note that this also works under Windows.)
    \end{description}
    \item 
        Equations, such as \begin{displaymath}
	    f(x)=\sum_{k\in\mathbb{Z}}c_{j_{0},k}\phi_{j_{0},k}(x)+\sum_{j=j_{0}}^{\infty}d_{j,k}\psi_{j,k}(x),
        \end{displaymath}
    where $\phi_{j,k}(x)=2^{-j/2}\phi(2^{j}x-k)$ and $\psi_{j,k}(x)=2^{-j/2}\psi(2^{j}x-k)$, for suitable functions $\psi(x)$ and $\phi(x)$.
\end{enumerate}
\end{document}

